% =====================================================================
%  7,320 Quantitative Risk Management (HSG, Autumn 2026) -- Assignment
%  Campanile, Giannattasio
%
%  Compile with pdfLaTeX (twice).  Files:
%     main.tex                 this file: preamble, title, abstract
%     sections/*.tex           one file per section (text is edited there)
%     generated/numbers.tex    every number quoted in the text  [written by scripts/08_build_report.py]
%     tables/*.tex             table bodies                     [idem]
%     figures/*.pdf            figures                          [idem]
%     roffin.cls               document class of the course template
% =====================================================================

% --- Compatibility of roffin.cls (2008) with a current LaTeX kernel -----------------------
% roffin.cls replaces \document and \enddocument by copies of the 2008 kernel code.  With a
% current kernel this breaks natbib (citations print as "[?]") and raises "Extra \endgroup".
% We keep the kernel's own definitions and restore them after the class is loaded.
\let\qrmkerneldocument\document
\let\qrmkernelenddocument\enddocument
\documentclass{roffin}
\makeatletter
\let\document\qrmkerneldocument
\let\enddocument\qrmkernelenddocument
% the class's \enddocument recorded the last page for the running head "first--last": do it here
\AtEndDocument{\immediate\write\@auxout{\string\endpage{\thepage}}}
% roffin.cls defines its own equation* environment, which clashes with amsmath: free the name
\expandafter\let\csname equation*\endcsname\relax
\expandafter\let\csname endequation*\endcsname\relax
% roffin.cls stores the full stop in \thesection ("1."), so that \ref prints "Section 1.".
% Keep the printed headings unchanged ("1. Title", "1.1 TITLE") but make references clean.
\renewcommand\thesection{\arabic{section}}
\renewcommand\thesubsection{\thesection.\arabic{subsection}}
\renewcommand\thesubsubsection{\thesubsection.\arabic{subsubsection}}
\def\qrm@dot@section{.}
\def\@seccntformat#1{\csname the#1\endcsname\@nameuse{qrm@dot@#1}\hskip .5em}
\makeatother

\usepackage{amsmath,amssymb,graphicx,booktabs}
\graphicspath{{figures/}}

% --- Notation -------------------------------------------------------------------------------
\newcommand{\E}{\mathbb{E}}
\renewcommand{\P}{\mathbb{P}}
\newcommand{\Var}{\mathrm{Var}}
\newcommand{\Cov}{\mathrm{Cov}}
\newcommand{\Corr}{\mathrm{Corr}}
\newcommand{\VaR}{\mathrm{VaR}}
\newcommand{\ES}{\mathrm{ES}}
\newcommand{\UL}{\mathrm{UL}}
\newcommand{\MDR}{\mathrm{MDR}}
\newcommand{\ind}{\mathbf{1}}
\newcommand{\dhat}{\hat d_k}

% --- Numbers from the code --------------------------------------------------------------------
% \val{key|format} prints the value defined in generated/numbers.tex; a value that has not been
% computed yet is printed as a bold [key], so that it cannot be overlooked.
\makeatletter
\newcommand{\qrmdef}[2]{\expandafter\gdef\csname qrm@\detokenize{#1}\endcsname{#2}}
\DeclareRobustCommand{\val}[1]{%
  \ifcsname qrm@\detokenize{#1}\endcsname\csname qrm@\detokenize{#1}\endcsname
  \else\mbox{\textbf{[\texttt{\detokenize{#1}}]}}\fi}
\makeatother
\InputIfFileExists{generated/numbers.tex}{}{}
% Visible reminder for text that can only be written once the CLIMADA results exist.
\newcommand{\todo}[1]{\textbf{[TO DO: #1]}}
% Boxed key message at the start of each section (management reading path).
\newcommand{\keymsg}[1]{\par\medskip\noindent\fbox{\parbox{\dimexpr\textwidth-2\fboxsep-2\fboxrule\relax}{\textbf{Key message.} #1}}\par\medskip}

% --- Tables: scale down to the text width only when a table is wider -------------------------
\newsavebox{\qrmbox}
\newcommand{\fittable}[1]{\sbox{\qrmbox}{#1}%
  \ifdim\wd\qrmbox>\textwidth\resizebox{\textwidth}{!}{\usebox{\qrmbox}}\else\usebox{\qrmbox}\fi}
% \qrmtable{name} and \qrmfigure{name} include tables/name.tex and figures/name.pdf; if the file
% has not been produced yet (question xi before the CLIMADA run) a visible note is printed instead.
\newcommand{\qrmtable}[1]{\IfFileExists{tables/#1.tex}%
  {{\small\setlength{\tabcolsep}{3.5pt}\fittable{\input{tables/#1}}}}%
  {\todo{table \texttt{\detokenize{#1}} is written by the scripts}}}
\newcommand{\qrmfigure}[1]{\IfFileExists{figures/#1.pdf}%
  {\includegraphics[width=\textwidth]{#1}}%
  {\todo{figure \texttt{\detokenize{#1}} is written by the scripts}}}

\begin{document}

\jname{Quantitative Risk Management Course}
\jyear{2026/27}
\jvol{}
\jnum{}
\doi{}
\accessdate{1 December 2026}
\copyrightinfo{Campanile, Giannattasio 2026.}

\markboth{Campanile, Giannattasio - University of St. Gallen (HSG)}{How Heavy Is the Tail?}

\title{How Heavy Is the Tail? Value-at-Risk and Expected Shortfall of a Credit Portfolio under Empirical, Gaussian and $t$-Copula Market Factors}

\author{Flavio Massimo Campanile$^1$, Pietro Tommaso Giannattasio$^2$}
\affil{$^1$University of St. Gallen (HSG), LUISS University}
\affil{$^2$University of St. Gallen (HSG)}

\maketitle

\begin{abstract}
We measure the risk of extreme loss of a portfolio of 100 corporate bonds over one week, at the $99.9\%$ level used by Basel regulation. Defaults cluster because the issuers' asset values move with the Swiss and US equity markets, so the result depends on how the two markets are assumed to move together. Holding every default probability fixed, the $99.9\%$ value-at-risk (VaR) is USD~\val{q8.M2.VaR_ref|mn2} million if market returns are Gaussian, USD~\val{q8.M3.VaR_ref|mn2} million if the two markets are allowed to crash together ($t$-copula), and USD~\val{q8.M1.VaR_ref|mn2} million if history is replayed as observed. The fat left tail of each market therefore matters far more than joint crashes. The Basel IRB capital, USD~\val{q10.K|usd}, covers only about a third of the model-based unexpected loss, because the portfolio is more exposed to the market than the regulatory correlations assume. Winter storms, assessed with the CLIMADA climate-risk platform, change expected loss and expected shortfall by amounts of the order of a few US dollars.
\end{abstract}

% =====================================================================
% 1. Executive summary
% Management version: every number comes from generated/numbers.tex through \val{...};
% results and conclusions are those of the full technical analysis (branch extra-features-version).
% =====================================================================
\section{Executive Summary}
\label{sec:summary}

\paragraph{The question.}
A single default rarely threatens a lender; many defaults at the same time do. In our portfolio of $100$ corporate bonds (USD~\val{q5.total_exposure|usd} of exposure) defaults cluster because the issuers' asset values move with two equity markets, the Swiss Performance Index (SPI) and the S\&P~500 (SPX). The regulatory loss measure, the $99.9\%$ value-at-risk (VaR) over the horizon of one week, is therefore driven by \emph{how the two markets are assumed to move together}. We hold every issuer's default probability fixed, so the expected loss is USD~\val{q5.EL|usd} in all cases, and compare three views of the market:
\begin{itemize}
\item \textbf{M1, history:} the observed weekly returns since 2000, replayed as they occurred;
\item \textbf{M2, Gaussian:} the textbook bell-curve model, in which joint crashes are very rare;
\item \textbf{M3, $t$-copula:} bell-curve markets that are allowed to crash \emph{together} more often.
\end{itemize}

\begin{table}[tb]
\caption{Key figures (USD; two million simulated scenarios per model). Unexpected loss $=$ VaR $-$ expected loss.}
\label{tab:key}
\centering
{\small\setlength{\tabcolsep}{4pt}
\begin{tabular}{@{}lrrr@{}}
\toprule
 & M2 Gaussian & M3 $t$-copula & M1 history \\
\midrule
VaR $99.9\%$ & \val{q8.M2.VaR_ref|usd} & \val{q8.M3.VaR_ref|usd} & \val{q8.M1.VaR_ref|usd} \\
\quad in \% of exposure & \val{q8.M2.VaR_pct_exposure|pct1} & \val{q8.M3.VaR_pct_exposure|pct1} & \val{q8.M1.VaR_pct_exposure|pct1} \\
Expected shortfall $99.9\%$ & \val{q8.M2.ES_ref|usd} & \val{q8.M3.ES_ref|usd} & \val{q8.M1.ES_ref|usd} \\
Unexpected loss & \val{q8.M2.UL_ref|usd} & \val{q8.M3.UL_ref|usd} & \val{q8.M1.UL_ref|usd} \\
\quad multiple of Basel IRB capital & \val{q10.ratio_UL_M2|f2} & \val{q10.ratio_UL_M3|f2} & \val{q10.ratio_UL_M1|f2} \\
\bottomrule
\end{tabular}}
\end{table}

\paragraph{Main findings.}
\begin{enumerate}
\item \textbf{Fat tails matter, joint crashes much less.} Allowing the two markets to crash together (M3) raises VaR by only \val{q8.M3.VaR_vs_M2|pct1}\% over the Gaussian model, although the data clearly prefer M3. Replaying history (M1) raises it by \val{q8.M1.VaR_vs_M2|pct0}\%, because each market on its own has a much fatter left tail than a bell curve (Table~\ref{tab:key}, Section~\ref{sec:results}).
\item \textbf{The historical figure rests on two weeks.} The M1 result is essentially the loss in the second-worst week since 2000; it would change if the sample were different.
\item \textbf{$10{,}000$ simulations are not enough to rank similar models.} At that size the VaR has a standard error of \val{derived.se_rel_min|pct0}--\val{derived.se_rel_max|pct0}\%, larger than the gap between M2 and M3. Our conclusions rest on two million scenarios.
\item \textbf{Through time, history is volatile and the Gaussian model is not.} On rolling two-year windows the Gaussian VaR hardly moves, while the historical VaR jumps by up to \val{q9.ratio_M1_M2_max|excess_pct0}\% above it whenever a crash week enters the window (Section~\ref{sec:rolling}).
\item \textbf{Regulatory capital is about a third of the model-based figure.} The Basel IRB capital is USD~\val{q10.K|usd} (risk-weighted assets USD~\val{q10.RWA|usd}), \val{q10.ratio_UL_M2|inv_pct0}\% of the Gaussian unexpected loss. The gap is due to the portfolio's high exposure to the market, not to the small number of names (Section~\ref{sec:irb}).
\item \textbf{Winter storms are immaterial for this portfolio.} Storm damage to the collateral, assessed with the CLIMADA platform, raises the expected loss by \val{q11.dEL_exact|ppm} parts per million and expected shortfall (ES) by about USD~\val{q11.first_order.M2.addon_tail_usd|usd2}. Even under a severe storm stress, or if the worst storm years coincided with the worst credit years, ES would move by less than $0.02\%$ (Section~\ref{sec:climate}).
\end{enumerate}

\paragraph{Implications for risk management.}
The choice of the \emph{marginal} distribution of market returns is the decision that drives tail risk; the dependence model is secondary. A tail estimate should always be reported with its simulation error, and differences between models smaller than that error should not be acted upon. For a portfolio with strong market exposure, the standard IRB capital is not conservative, and the gap belongs to the supervisory review of Pillar~2. Under the present climate, Swiss winter storms do not change the capital picture of this portfolio.

\paragraph{How to read this report.}
Sections~\ref{sec:setting}--\ref{sec:conclusion} give the results in non-technical terms, each opening with its key message. The derivations required by the assignment are collected in the Technical Appendix, and Table~\ref{tab:map} at its end maps every question of the assignment to the place where it is answered. All figures are produced by the attached code, which reproduces every number in this report.

% =====================================================================
% 2. Portfolio, model and data -- questions (i)-(iv); derivations in Appendix A-B
% =====================================================================
\section{Portfolio, Model and Data}
\label{sec:setting}

\keymsg{Defaults cluster only through the two equity markets. Both markets have a much fatter left tail than a bell curve, and they crash together more often than a Gaussian model allows; the three models of the market are built to separate these two effects.}

\subsection{THE PORTFOLIO}
The portfolio consists of $100$ corporate bonds of firms with branches in Switzerland and the United States. Total exposure is USD~\val{q5.total_exposure|usd}, of which \val{q5.share_exposure_inv_grade|pct0}\% is rated BBB or better. One-year default probabilities $\pi_k$ range from $0.01\%$ (AAA) to $20\%$ (CCC) and recovery rates $R_k$ from \val{q5.R_min|f2} to \val{q5.R_max|f2}. The expected loss is USD~\val{q5.EL|usd}; \val{q5.EL_share_B_CCC|pct0}\% of it comes from the B and CCC names. The issuers are located in \val{q5.n_cantons|int} Swiss cantons, which matters only for the storm analysis of Section~\ref{sec:climate}.

\subsection{HOW DEFAULTS ARISE}
We use the threshold model of the assignment, a version of Merton's model \citep[Ch.~2]{DeGiorgi26}: firm $k$ defaults when the return on its assets over the week, $Y_k$, falls below a threshold $d_k$, i.e.\ $I_k=\ind_{\{Y_k\le d_k\}}$. The asset return combines the two market returns $\Theta=(\Theta_1,\Theta_2)'$ (SPI and SPX) with a firm-specific shock:
\begin{equation}
\label{eq:Yk}
Y_k=\sqrt{\lambda_k}\,a_k'\Theta+\sqrt{1-\lambda_k}\,s_k\,\epsilon_k ,
\end{equation}
where the shocks $\epsilon_k\sim N(0,1)$ are i.i.d.\ and independent of $\Theta$. The loadings $a_k$ say how much each firm depends on the Swiss and the US market, and $\lambda_k$ is the share of its risk that comes from the market (on average \val{q5.lam_mean|f2}). Three results, derived in Appendix~A, organise the analysis:
\begin{itemize}
\item \emph{Given the market, defaults are independent coin flips} (question i). Firm $k$ defaults with probability
\[
p_k(\theta)=\Phi\!\left(\frac{d_k-\sqrt{\lambda_k}\,a_k'\theta}{\sqrt{1-\lambda_k}\,s_k}\right),\qquad I_k\mid\Theta\sim\mathrm{Bernoulli}\big(p_k(\Theta)\big).
\]
A market fall raises every $p_k$ at once: this is the only channel through which defaults cluster.
\item \emph{The loss} (question ii) is $L=\sum_kE_k(1-R_k)I_k$, the exposure net of recovery of the names that default; its mean, the expected loss, depends only on the $\pi_k$.
\item \emph{The scale of $Y_k$} (question iii) is $s_k=\sqrt{a_k'\Sigma a_k}$, with $\Sigma$ the covariance matrix of $\Theta$.
\end{itemize}
Each threshold $d_k$ is set so that the firm defaults with its given probability, $\P[Y_k\le d_k]=\pi_k$. As a result, the three market models below produce the same expected loss, and any difference in VaR or ES comes from the \emph{shape} of the market distribution alone (Appendix~A).

\subsection{MARKET DATA}
We use weekly (Wednesday to Wednesday, non-overlapping) log-returns of the two indices from \val{q4.first_week|date} to \val{q4.last_week|date}, $n=\val{q4.n_weeks|int}$ weeks, on days when both markets traded. Both series include dividends: the SPX file is the total-return index. The data show two departures from a bell curve:
\begin{itemize}
\item \emph{Fat left tails.} Both markets are left-skewed and strongly leptokurtic: the Jarque--Bera test rejects normality far beyond any conventional level. The SPX fell by \val{derived.cel1_spx|pct1}\% in the week ending \val{derived.cel1_week|date}, a seven-standard-deviation event for a Gaussian model with the estimated volatility.
\item \emph{Joint crashes.} When the SPI is in its worst $1\%$ of weeks, the SPX is in its own worst $1\%$ in \val{q4.joint_exceed_emp.0.01|pct0}\% of cases, against \val{q4.joint_exceed_gauss.0.01|pct0}\% under a Gaussian model with the same correlation (\val{q4.corr_all|f2}).
\end{itemize}

\subsection{THREE VIEWS OF THE MARKET}
The three models are combinations of \emph{margins} (how risky each market is on its own) and a \emph{copula} (how the two move together), which makes them comparable:
\begin{description}
\item[M1] the empirical distribution: each historical week is a scenario with probability $1/n$; nothing is estimated.
\item[M2] a bivariate Gaussian distribution: Gaussian margins, Gaussian copula.
\item[M3] Gaussian margins with a $t$-copula, which allows joint crashes.
\end{description}
M3 and M2 differ only in the copula, so their gap measures the effect of joint crashes; M1 and M2 differ in both, so their gap measures the full cost of the Gaussian assumption. M2 and M3 are estimated by maximum likelihood (question iv; derivations in Appendix~B). Table~\ref{tab:mle} reports the estimates. The $t$-copula has $\hat\nu=\val{q4.M3.params.5|f1}$ degrees of freedom and a likelihood-ratio statistic of \val{q4.LR|f1} against the Gaussian copula ($p$-value below $10^{-9}$): \emph{joint crashes are statistically real}. But M3 keeps bell-curve margins, which is where the data depart most from the Gaussian model.

\begin{table}[tb]
\caption{Maximum-likelihood estimates of M2 and M3, weekly log-returns (standard errors in parentheses).}
\label{tab:mle}
\centering
\qrmtable{tab_mle}
\end{table}

% =====================================================================
% 3. Risk measures and simulation -- questions (v)-(vii); details in Appendix C-D
% =====================================================================
\section{How We Measure Tail Risk}
\label{sec:method}

\keymsg{A $99.9\%$ risk measure is set by the ten worst of $10{,}000$ scenarios, so its simulation error must be measured. We do so with a reference run of two million scenarios.}

\paragraph{Risk measures.}
Following Chapter~1 of the lecture notes \citep{DeGiorgi26}, the value-at-risk $\VaR_\alpha$ is the loss that is exceeded with probability at most $1-\alpha$, and the expected shortfall $\ES_\alpha$ is the average loss in that worst $1-\alpha$ of cases. ES is coherent, while VaR is not subadditive in general. The unexpected loss $\UL_\alpha=\VaR_\alpha-\E[L]$ is the amount that capital is meant to cover. We use $\alpha=99.9\%$, the Basel IRB level. The formal definitions and their estimators for a simulated, discrete loss are given in Appendix~D.

\paragraph{Simulation (questions v--vii).}
Each scenario draws the two market returns from the model at hand (a historical week under M1, a Gaussian draw under M2, a $t$-copula draw with Gaussian margins under M3; algorithms in Appendix~C) and $100$ independent firm-specific shocks. It then computes the asset returns $Y_k$ by (\ref{eq:Yk}), the defaults $I_k$ and the portfolio loss $L$. As required, we simulate $10{,}000$ scenarios per model. The three models use the \emph{same} random numbers, so that differences between them are not blurred by different draws.

\paragraph{Default thresholds.}
The threshold $d_k$ must satisfy $\P[Y_k\le d_k]=\pi_k$. Following the assignment's hint, we take the empirical $\pi_k$-quantile of the $10{,}000$ simulated $Y_k$: each firm then defaults in exactly a share $\pi_k$ of the scenarios, and the average simulated loss equals the expected loss to the cent. For the best ratings this quantile rests on very few scenarios (one for AAA), so we also compute a \emph{reference threshold} that solves the same condition on the model itself, without simulation error (Appendix~C).

\paragraph{Measuring precision.}
A $99.9\%$ VaR from $10{,}000$ scenarios is the tenth-largest loss, a thin basis for an estimate. We therefore also run every model on two million scenarios, organised as $200$ independent repetitions of the $10{,}000$-scenario experiment. The spread across repetitions gives the standard error of each estimate and of each difference between models.

\begin{figure}[tb]
\centering
\qrmfigure{fig_loss_hist}
\caption{Simulated distribution of the portfolio loss under the three models (question vii), $10{,}000$ scenarios each, logarithmic vertical axis. The models agree in the body and separate only in the far right tail.}
\label{fig:hist}
\end{figure}

Figure~\ref{fig:hist} shows the three simulated loss distributions. In most scenarios no or few firms default; the models are indistinguishable in this region and differ only in the far tail, where the $99.9\%$ measures are determined.

% =====================================================================
% 4. VaR and ES -- question (viii)
% =====================================================================
\section{Value-at-Risk and Expected Shortfall}
\label{sec:results}

\keymsg{Historical returns give a VaR \val{q8.M1.VaR_vs_M2|pct0}\% above the Gaussian model; allowing joint crashes adds only \val{q8.M3.VaR_vs_M2|pct1}\%. The fat tails of each market, not their dependence, drive tail risk.}

\begin{table}[tb]
\caption{Question (viii), baseline run: $99.9\%$ VaR and ES (USD) from $10{,}000$ scenarios, with the empirical-quantile thresholds $\hat d_k$ and with the reference thresholds $d_k$. Standard errors from 200 repetitions; last row: the difference M3 $-$ M2. The expected loss is USD~\val{q5.EL|usd} in all models.}
\label{tab:var}
\centering
\qrmtable{tab_var_es}
\end{table}

\begin{table}[tb]
\caption{Reference run: VaR, ES and unexpected loss (UL $=$ VaR $-$ EL), USD, two million scenarios. ``\% exp.'': VaR in per cent of total exposure. Standard errors of the M3 $-$ M2 difference in parentheses.}
\label{tab:varref}
\centering
\qrmtable{tab_var_es_ref}
\end{table}

\paragraph{The numbers.}
With $10{,}000$ scenarios (Table~\ref{tab:var}) the $99.9\%$ VaR is USD~\val{q8.M2.VaR_10k_quantile_thr|usd} under M2, USD~\val{q8.M3.VaR_10k_quantile_thr|usd} under M3 and USD~\val{q8.M1.VaR_10k_quantile_thr|usd} under M1. With two million scenarios (Table~\ref{tab:varref}) the same figures are USD~\val{q8.M2.VaR_ref|usd}, USD~\val{q8.M3.VaR_ref|usd} and USD~\val{q8.M1.VaR_ref|usd}, i.e.\ \val{q8.M2.VaR_pct_exposure|pct1}\%, \val{q8.M3.VaR_pct_exposure|pct1}\% and \val{q8.M1.VaR_pct_exposure|pct1}\% of exposure. ES exceeds VaR by \val{derived.ES_over_VaR_min|pct0}--\val{derived.ES_over_VaR_max|pct0}\% and ranks the models in the same way.

\paragraph{Joint crashes (M3 against M2) matter little.}
The $t$-copula raises VaR by \val{q8.M3.VaR_vs_M2|pct1}\% (USD~\val{q8.M3_minus_M2.dVaR_ref|usd}, standard error about USD~\val{q8.M3_minus_M2.se_dVaR_ref|usd}) and ES by \val{q8.M3.ES_vs_M2|pct1}\%: statistically detectable, economically minor. Three reasons explain it. Each firm sees the markets only through a weighted sum of the two returns, and summing two bell-shaped, already highly correlated returns produces an almost bell-shaped result whatever their copula. The estimated copula ($\hat\nu\approx15$) is only a moderate departure from the Gaussian one. And the firm-specific shocks remain Gaussian and independent, which dilutes the market's tail behaviour.

\paragraph{Fat tails (M1 against M2) matter a lot, but rest on few weeks.}
Replaying history raises VaR by \val{q8.M1.VaR_vs_M2|pct0}\% and ES by \val{q8.M1.ES_vs_M2|pct0}\%; the $99.9\%$ quantile of the number of defaults rises from \val{q8.M2.defaults_q999|int} to \val{q8.M1.defaults_q999|int}. At the $99\%$ level the gap is only \val{derived.M1_vs_M2_99|pct0}\%: the models separate in the far tail. Each week has probability $1/n=\val{derived.one_over_n|pct3}\%$ under M1, so the $0.1\%$ tail consists of one or two weeks: those ending \val{derived.cel1_week|date} (expected loss given that week USD~\val{derived.cel1|mn2} million) and \val{derived.cel2_week|date} (USD~\val{derived.cel2|mn2} million). Falls of this size are practically impossible under Gaussian margins. The M1 figure is therefore realistic about crashes but fragile: it would change with a different sample, and it cannot produce a week worse than the worst one observed.

\paragraph{Precision: $10{,}000$ scenarios cannot rank M2 and M3.}
With $10{,}000$ scenarios the standard error of VaR is \val{derived.se_rel_min|pct0}--\val{derived.se_rel_max|pct0}\% of its value (Table~\ref{tab:var}). The empirical-quantile thresholds give the more precise estimate, because recalibrating in each sample offsets part of the sampling noise. The noise in the M3 $-$ M2 difference is \val{derived.noise_ratio|f1} times the difference itself, and M3 comes out \emph{below} M2 in \val{q8.M3_minus_M2.ref_thr.p_M3_below_M2_crn|pct0}\% of the repetitions; our own baseline run shows opposite signs with the two thresholds. Statements on the size of the copula effect rest on the two-million-scenario run alone.

\paragraph{Name concentration is minor.}
Removing all firm-specific risk, i.e.\ moving to an infinitely granular portfolio with the same market exposure, lowers VaR by at most \val{derived.gran_share_VaR_max|pct1}\%: with $100$ names, tail risk is market risk.

% =====================================================================
% 5. Dynamics -- question (ix)
% =====================================================================
\section{Risk Through Time}
\label{sec:rolling}

\keymsg{Re-estimated every week on the last two years of data, the Gaussian VaR is almost constant, the $t$-copula VaR is indistinguishable from it, and the historical VaR swings between USD~\val{q9.min.VaR_M1|mn2} and \val{q9.max.VaR_M1|mn2} million depending on whether a crash week is in the window.}

\paragraph{Design (question ix).}
At each date we re-estimate the three models on the last $500$ daily returns (about two years), recompute $s_k$ and the thresholds so that every firm keeps its default probability $\pi_k$, and simulate $50{,}000$ scenarios. Two years contain only $100$ non-overlapping weeks, too few for a $t$-copula or an empirical $99.9\%$ quantile, so within each window we use the $496$ overlapping five-day returns. The window moves forward by one week, giving \val{q9.n_windows|int} dates from \val{q9.first|date} to \val{q9.last|date}. To judge whether an M3 $-$ M2 difference is real, we compare it with a \emph{placebo band}: the spread of the same difference when the two models are identical, which measures pure simulation noise (standard deviation \val{q9.placebo.VaR_sd|pct1}\% for VaR).

\begin{figure}[p]
\centering
\qrmfigure{fig_rolling}
\caption{Rolling estimates on windows of 500 daily returns. Top two panels: $99.9\%$ VaR and ES with thresholds recalibrated to $\pi_k$ in every window. Third panel: VaR with thresholds frozen at their full-sample values. Lower panels: dependence parameters and degrees of freedom of the $t$-copula.}
\label{fig:rolling}
\end{figure}

\begin{table}[tb]
\caption{Rolling $99.9\%$ VaR and ES (USD), \val{q9.n_windows|int} windows of 500 daily returns.}
\label{tab:rolling}
\centering
\qrmtable{tab_rolling}
\end{table}

\paragraph{Gaussian model: almost constant.}
The M2 VaR stays between USD~\val{q9.min.VaR_M2|usd} and USD~\val{q9.max.VaR_M2|usd} (Figure~\ref{fig:rolling}, Table~\ref{tab:rolling}), although market volatility varies by a factor of \val{q9.sigma_ratio_max_min|f1} across windows. Because the thresholds are reset to the same default probabilities, higher volatility is absorbed by lower thresholds; what remains is the correlation between the markets (correlation between rolling VaR and rolling $\hat\rho$: \val{q9.corr_VaRM2_rho|f2}).

\paragraph{$t$-copula: within the noise for VaR, slightly above it for ES.}
The M3 $-$ M2 difference in VaR averages \val{q9.dVaR_mean|spct1}\% and lies inside the placebo band in \val{q9.share_dVaR_in_band|pct0}\% of the windows. Two-year windows are too short to pin down the degrees of freedom ($\hat\nu$ is at its upper bound in \val{q9.share_nu_at_bound|pct0}\% of them). For ES the difference averages \val{q9.dES_mean|spct1}\% and exceeds the band in \val{q9.share_dES_above_band|pct0}\% of the windows, mostly in crisis windows with strong tail dependence: joint crashes show up in ES before VaR, but by less than 10\%.

\paragraph{Historical model: driven by crash weeks.}
The M1 VaR ranges from USD~\val{q9.min.VaR_M1|usd} to USD~\val{q9.max.VaR_M1|usd} (window ending \val{q9.argmax_VaR_M1|date}), with a standard deviation \val{q9.ratio_std_M1_M2|f0} times that of M2. With $496$ observations, the $99.9\%$ tail of the market is the single worst week in the window: the VaR jumps when a crash enters (2008, 2011, March 2020) and drops 500 days later when it leaves, the ``ghost effect'' of historical simulation. In calm windows M1 is even below M2 (\val{q9.share_M1_below_M2|pct0}\% of the windows).

\paragraph{What recalibration hides.}
Resetting the thresholds makes the model ``through-the-cycle'', which fits rating-based default probabilities like ours. If instead the thresholds are frozen at their full-sample values (third panel), higher volatility raises default probabilities: the Gaussian VaR then ranges from USD~\val{q9.min.VaRfix_M2|usd} to USD~\val{q9.max.VaRfix_M2|usd}, and the common volatility cycle dominates the differences between models.

% =====================================================================
% 6. Basel IRB -- question (x)
% =====================================================================
\section{Regulatory Capital: the Basel IRB Formula}
\label{sec:irb}

\keymsg{The IRB capital is USD~\val{q10.K|usd} (RWA USD~\val{q10.RWA|usd}), about a third of the model-based unexpected loss. Almost the whole gap comes from one assumption: the portfolio's firms are far more exposed to the market than the regulatory correlations allow for.}

\subsection{CAPITAL CHARGE BY COUNTERPARTY}
For each counterparty we apply the IRB formula of the assignment \citep{BCBS05} with $\mathrm{AVC}=1$ and maturity $M=1$:
\begin{align*}
\rho_k&=0.12\,w_k+0.24\,(1-w_k),\qquad w_k=\frac{1-e^{-50\pi_k}}{1-e^{-50}},\\
K_k&=(1-R_k)\left[\Phi\!\left(\frac{\Phi^{-1}(\pi_k)+\sqrt{\rho_k}\,\Phi^{-1}(0.999)}{\sqrt{1-\rho_k}}\right)-\pi_k\right]\mathrm{MA}(\pi_k,M).
\end{align*}
The term in brackets is the default probability in a $1$-in-$1{,}000$ bad year minus the expected one, which provisions already cover. Table~\ref{tab:irb} reports the results by rating class: $\rho_k$ falls from $0.24$ for the best ratings to $0.12$ for CCC, and $K_k$ rises from \val{q10.Kbar_pct.AAA|f2}\% of exposure for AAA to \val{q10.Kbar_pct.CCC|f2}\% for CCC. Summing, $K=\sum_kE_kK_k=$ USD~\val{q10.K|usd} (\val{q10.K_pct|pct2}\% of exposure) and $\mathrm{RWA}=12.5\,K=$ USD~\val{q10.RWA|usd}; B and CCC names account for \val{derived.K_share_B_CCC|pct0}\% of $K$.

\begin{table}[tb]
\caption{Basel IRB quantities by rating class ($\mathrm{AVC}=1$, $M=1$). $\bar K_k$, $\overline{\mathrm{LGD}}$: averages within the class; amounts in USD. Last column: average market share of risk $\lambda_k$ in our model.}
\label{tab:irb}
\centering
\qrmtable{tab_irb}
\end{table}

\paragraph{What $\mathrm{MA}(\pi_k,1)$ means.}
With $M=1$ the maturity adjustment is neutral for every counterparty:
\[
\mathrm{MA}(\pi_k,1)=\frac{1+(1-2.5)\,b(\pi_k)}{1-1.5\,b(\pi_k)}=1 .
\]
The core formula covers default within the one-year horizon. A longer loan can also lose value within the year through a downgrade, and the adjustment adds capital for this migration risk. With a maturity of exactly one year nothing of the loan is left after the horizon to be downgraded, so $K_k$ is a pure default charge.

\subsection{WHY IRB CAPITAL AND MODEL-BASED LOSS DIFFER}
The unexpected loss is USD~\val{q8.M2.UL_ref|usd} under M2 and USD~\val{q8.M3.UL_ref|usd} under M3, \val{q10.ratio_UL_M2|f2} and \val{q10.ratio_UL_M3|f2} times $K$. The IRB formula is itself the $99.9\%$ unexpected loss of a specific model: one Gaussian factor, the prescribed correlation $\rho_k$, and infinitely many small loans \citep{Vasicek02,Gordy03}. Table~\ref{tab:gap} moves from that model to ours one assumption at a time.

\begin{table}[tb]
\caption{From the IRB capital to the model-based unexpected loss at $99.9\%$ (reference run), one assumption at a time.}
\label{tab:gap}
\centering
\qrmtable{tab_irb_gap}
\end{table}

\begin{itemize}
\item \emph{Correlation level: almost the entire gap.} In our model $\lambda_k$ plays the role of $\rho_k$ (Appendix~A). The portfolio's $\lambda_k$ average \val{q10.lambda_mean|f2}, against regulatory values of $0.12$--$0.24$, and do not fall with the default probability. Using $\rho_k=\lambda_k$ in the IRB formula alone multiplies the capital by \val{derived.K_lambda_ratio|f2}.
\item \emph{Two factors instead of one:} small, because the two indices are highly correlated and all loadings are positive.
\item \emph{Granularity:} $100$ names instead of infinitely many adds \val{derived.gran_share_UL_M2|pct0}\% of the unexpected loss under M2.
\item \emph{Shape of the factors:} the $t$-copula adds \val{derived.UL_M3_vs_M2|pct1}\%, historical returns \val{derived.UL_M1_vs_M2|pct0}\%.
\item \emph{Horizon:} immaterial under M2, since after calibrating $d_k$ only the shape of the factor distribution matters (Appendix~A).
\end{itemize}

\subsection{WHY A REGULATOR PREFERS THE SIMPLER FORMULA}
\begin{itemize}
\item \emph{Portfolio invariance:} the capital for a loan depends only on that loan, so capital is additive and translates directly into risk weights \citep{Gordy03}.
\item \emph{Comparability:} two banks holding the same loan hold the same capital. Model-based capital would differ for reasons unrelated to risk; in our case, the choice between M1, M2 and M3 alone moves the unexpected loss by about \val{derived.UL_M1_vs_M2|pct0}\%.
\item \emph{Less room for model risk and arbitrage:} correlations and tail shapes, the inputs to which the tail is most sensitive and which are hardest to backtest at $99.9\%$, are fixed by the regulator.
\item \emph{Verifiability:} a closed formula can be checked line by line; a simulated figure carries Monte Carlo error (Section~\ref{sec:results}).
\end{itemize}
The price is that the formula ignores what is specific to a portfolio, such as concentration and correlations above the prescribed ones; Basel leaves these to the supervisory review of Pillar~2. Our portfolio shows that this omission need not be conservative.

% =====================================================================
% 7. Physical climate risk -- question (xi); formulas in Appendix E
% =====================================================================
\section{Physical Climate Risk: Winter Storms}
\label{sec:climate}

\keymsg{Swiss winter storms barely damage the collateral: the expected loss rises by \val{q11.dEL_exact|ppm} parts per million and ES by about USD~\val{q11.first_order.M2.addon_tail_usd|usd2}. A more diversified portfolio would not change this; only a strong link between storm years and credit years, or a far more destructive hazard, would make it visible, and even then by less than $0.02\%$ of ES.}

\paragraph{What kind of risk, and what kind of statement.}
Climate risks are either \emph{transition} risks (regulation, technology, preferences in the move to a low-carbon economy) or \emph{physical} risks, acute or chronic \citep{Acharya23}. Winter storms are an acute physical risk; here they reach credit risk through the collateral, i.e.\ through the loss given default. Unlike supervisory stress tests, which report losses conditional on scenarios such as the NGFS pathways \citep{NGFS24}, our storm catalogue describes the \emph{present} climate probabilistically, so the overlay yields unconditional VaR and ES.

\paragraph{Hazard, exposure, vulnerability.}
Following the CLIMADA framework \citep{AznarSiguan19,HauserEtAl26}, the \emph{hazard} is CLIMADA's probabilistic set of European winter storms for Switzerland: \val{q11.hazard.n_events|int} events (historical storms and perturbed versions), each with an annual frequency $f_e$ and a gust-speed footprint. The \emph{exposure} is the recovery value $E_kR_k$ of each counterparty, placed at its coordinates (USD~\val{q5.recovery_value|usd} in total). The \emph{vulnerability} is the impact function of \citet{Welker21} for Swiss buildings, which turns gust speed into a mean damage ratio $\MDR_k(e)$. We evaluate the hazard at each counterparty's coordinates, the most precise information available; evaluating it at the centroid of the canton, as worded in the assignment, gives the same conclusions (Table~\ref{tab:climvar}).

\begin{figure}[tb]
\centering
\qrmfigure{fig_climate_hazard}
\caption{Left: gust footprint of the most damaging storm of the catalogue (Lothar, 26 December 1999) and location of the counterparties (circle area proportional to exposure). Right: impact function (mean damage ratio, logarithmic scale) and frequency-weighted distribution of the gust speeds at the counterparties.}
\label{fig:climhaz}
\end{figure}

\paragraph{The collateral-sensitivity parameter $\phi_k$.}
The climate-adjusted recovery rate is $R_k(e)=R_k(1-\phi_k\MDR_k(e))$.
\begin{itemize}
\item $\phi_k=0$: none of the recovery value is physically at the exposed site (financial collateral, guarantees, assets elsewhere). Storms cannot reduce recovery, and the counterparty is immune to the overlay.
\item $\phi_k=1$: the whole recovery value consists of physical assets at the site. Every unit of damage reduces recovery proportionally; a storm that destroys the site leaves nothing to recover.
\end{itemize}
We use $\phi_k=0.5$ as required. Two real-world features argue that the effective value is lower. In Switzerland, insurance of buildings against natural hazards, storms included, is mandatory, and in 19 cantons it is provided by the cantonal building insurer \citep{BAFU20}: insured damage does not reduce the bank's recovery. And the impact function, calibrated on the claims of the Zurich cantonal building insurer, already reflects today's building standards.

\paragraph{Design.}
On top of the $10{,}000$ credit scenarios of Section~\ref{sec:results}, and independently of them as the assignment prescribes, we draw $N_e\sim\mathrm{Poisson}(f_e)$ for every event; if several events hit a counterparty in the same year we keep the worst one, $m_k$. Defaults are unchanged; each defaulted name loses in addition the fraction $\phi\,m_k$ of what it would have recovered, so the loss becomes $L+\phi A$ with $A=\sum_kE_kR_km_kI_k$. We repeat the exercise on the two-million-scenario reference run.

\begin{table}[tb]
\caption{Question (xi): $99.9\%$ VaR and ES without the storm overlay (USD) and their increase with it, in USD and in parts per million (ppm; $1$~ppm $=0.0001\%$), $\phi_k=0.5$.}
\label{tab:clim}
\centering
\qrmtable{tab_climate}
\end{table}

\paragraph{Results against question (viii).}
Table~\ref{tab:clim} compares VaR and ES with and without storms. The increases are a few US dollars at most: no relative change exceeds \val{q11.assign.M2.dES|ppm}~ppm in the $10{,}000$-scenario run. The reason lies in the hazard. Gusts at the counterparties rarely exceed the $20$~m/s at which damage starts (Figure~\ref{fig:climhaz}); the largest damage ratio in the whole catalogue is \val{q11.hazard.mdr_max|pct2}\%, during Lothar, and the expected worst damage of a year is \val{q11.hazard.exp_worst_mean|ppm}~ppm of the collateral on average. The VaR changes are not comparable across models: with $10{,}000$ scenarios VaR is a single scenario, and its change is the damage of the one storm year paired with it. ES, an average over ten scenarios, is the reliable measure; its increase matches a closed-form first-order approximation (Appendix~E). Because storms and defaults are independent, the tail is hit no harder than the mean: ES rises by about \val{q11.first_order.M2.dES|ppm}~ppm against \val{q11.dEL_exact|ppm}~ppm for the expected loss.

\begin{table}[tb]
\caption{Reference run ($N=2{,}000{,}000$): increase of the expected loss (M2), VaR and ES due to the storm overlay (ppm), for the baseline and five variants. First row: first-order approximation. ES is the reliable measure; the VaR columns are dominated by a single scenario.}
\label{tab:climvar}
\centering
\qrmtable{tab_climate_variants}
\end{table}

\paragraph{A more geographically diversified portfolio.}
All counterparties are in Switzerland, smaller than a typical storm footprint: the most damaging event hits \val{q11.hazard.share_names_hit_by_top_event|pct0}\% of them at once. To mimic full diversification we let every counterparty draw its own, independent storm year (Table~\ref{tab:climvar}, ``Independent storm years''). ES then changes by \val{q11.ref.diversified.M2.dES|ppm}~ppm under M2, against \val{q11.ref.base.M2.dES|ppm}~ppm in the baseline, with no systematic difference across models. Under independence what matters is the \emph{average} hazard at the locations of the defaulting names, not its concentration. A diversified portfolio would thus show a different effect only if its locations were on average more exposed (North Sea coasts, cyclone or flood zones) or less exposed (south of the Alps) than Switzerland.

\paragraph{When concentration does matter: dependence.}
If the worst storm years were paired with the worst credit years (``Comonotone'' in Table~\ref{tab:climvar}), ES under M2 would rise by \val{q11.ref.comonotone.M2.dES|ppm}~ppm (USD~\val{q11.ref.comonotone.M2.dES_usd|usd}), \val{q11.ref.comonotone.M2.dES_vs_base|f0} times the baseline. Only then is concentration costly, since one bad storm year hits all Swiss names together. Even this extreme pairing moves ES by less than $0.02\%$; it bounds the cost of the independence assumption, the strongest simplification of the exercise.

\paragraph{Beyond the assignment: link to the course model and a hazard stress.}
\begin{table}[tb]
\caption{Storms through both credit channels (reference run, $\phi_k=0.5$, asset volatility $\sigma_A=25\%$ unless stated). $\Delta$EL exact; $\Delta$ES first-order. Lower block: all gust speeds raised by the stated percentage.}
\label{tab:climchannels}
\centering
\qrmtable{tab_climate_channels}
\end{table}
Our overlay is exactly the LGD channel of the credit model with physical risk of \citet{Pozdyshev25} presented in the course, with $1-e^{-\alpha}=\phi\MDR$ (Appendix~E). In their model the same storm also lowers the firm's asset value and hence raises its default probability. Adding this \emph{PD channel} (Table~\ref{tab:climchannels}) raises the expected-loss effect to \val{q11x.ref.both.dEL_exact|ppm}~ppm, \val{q11x.pd_share_EL|pct0}\% of it from the PD channel: a small jump in asset value moves a small default probability proportionally more than it moves the recovery. ES rises by USD~\val{q11x.ref.both.M2.dES_usd|usd2} under M2, and the IRB capital by \val{q11x.irb.s25.d_both|ppm}~ppm. Finally, a stylised high-hazard scenario that raises every gust speed by $20\%$ multiplies the expected damage by \val{q11x.stress.s020.Em_ratio_to_base|f1}, because the impact function is steeply convex; ES then rises by USD~\val{q11x.ref.g20.M2.dES_usd|usd0} under M2, still below $0.01\%$. For a Swiss portfolio the conclusion that winter storms are immaterial is robust to both extensions.

% =====================================================================
% 8. Conclusion
% =====================================================================
\section{Conclusion}
\label{sec:conclusion}

Holding default probabilities fixed, the $99.9\%$ VaR of the portfolio ranges from USD~\val{q8.M2.VaR_ref|mn2} million under Gaussian factors to USD~\val{q8.M1.VaR_ref|mn2} million under the empirical distribution, with the $t$-copula model only \val{q8.M3.VaR_vs_M2|pct1}\% above the Gaussian one. The tail of the loss is sensitive to the distribution of the market factors, but to their margins far more than to their copula: allowing the two indices to crash together matters little as long as each is Gaussian and the firm-specific shocks are Gaussian and independent. A natural extension would combine a $t$-copula with fat-tailed margins, obtaining the tail thickness of M1 without tying the $99.9\%$ quantile to two historical weeks.

The same message recurs elsewhere. Through time, the parametric models are stable because the thresholds are recalibrated, while the empirical model inherits the ghost effects of historical simulation. Against the regulatory benchmark, the model-based unexpected loss is \val{q10.ratio_UL_M2|f1} times the IRB capital, a matter of correlation levels, not of granularity or tail shape. The storm overlay, which acts on recoveries and is independent of the credit cycle by assumption, raises the expected loss by \val{q11.dEL_exact|ppm} parts per million and ES by about USD~\val{q11.first_order.M2.addon_tail_usd|usd2}, four orders of magnitude below the Monte Carlo error of the estimates. Concentration is irrelevant under independence and costs less than $0.02\%$ of ES even under dependence; the PD channel of \citet{Pozdyshev25} and a $20\%$ gust stress raise the effect several- and tenfold, yet ES moves by less than $0.01\%$.

\paragraph{Caveats.}
The empirical model depends on a handful of observations. With $10{,}000$ simulations the Monte Carlo error at $99.9\%$ is of the same order as some of the differences of interest, which is why our conclusions rest on two million scenarios. The comparison with the IRB formula is driven by the systematic weights $\lambda_k$ supplied with the synthetic portfolio; with weights in the regulatory range, model-based and regulatory capital would be much closer. Finally, the climate overlay covers one acute physical hazard under the present climate and through the collateral only: transition and chronic risks, liquidity and operational channels, and scenario-conditional losses in the sense of supervisory climate stress tests are outside its scope.

% =====================================================================
% Technical appendix: the derivations required by the assignment (iii of the conditions)
% =====================================================================
\appendix
\section*{Technical Appendix}

\subsection*{A. THE CREDIT MODEL: QUESTIONS (i)--(iii)}
\paragraph{Link to Merton's model.}
In Merton's model \citep[Ch.~2]{DeGiorgi26} the firm defaults if $V_{k,T}\le B_k$, i.e.\ if $Y_k=\log(V_{k,T}/V_{k,0})\le d_k=\log(B_k/V_{k,0})$. Model (\ref{eq:Yk}) keeps this structure, adds factor dependence \citep[Sect.~8.3]{McNeil05} and calibrates $d_k$ to $\pi_k$.

\paragraph{(i) Conditional default probability.}
Fix $\Theta=\theta$. Since $\sqrt{1-\lambda_k}\,s_k>0$, $Y_k\le d_k\iff\epsilon_k\le(d_k-\sqrt{\lambda_k}a_k'\theta)/(\sqrt{1-\lambda_k}\,s_k)$. As $\epsilon_k\sim N(0,1)$ is independent of $\Theta$, conditioning does not change its law, hence
\begin{equation}
\label{eq:condpd}
p_k(\theta):=\P[I_k=1\mid\Theta=\theta]=\Phi\!\left(\frac{d_k-\sqrt{\lambda_k}\,a_k'\theta}{\sqrt{1-\lambda_k}\,s_k}\right),
\end{equation}
that is, $I_k\mid\Theta\sim\mathrm{Bernoulli}(p_k(\Theta))$.
Given $\Theta$, the $I_k$ are functions of the independent $\epsilon_k$ only, so they are conditionally independent (a Bernoulli mixture model, \citealp[Sect.~8.4]{McNeil05}). By iterated expectations $\pi_k=\P[Y_k\le d_k]=\E[p_k(\Theta)]$, and joint defaults have probability $\E[p_j(\Theta)p_k(\Theta)]$: default correlation arises only through $\Theta$. With $a_k\ge0$ for all names, a fall of both indices raises every $p_k$.

\paragraph{(ii) Portfolio loss.}
A default costs the exposure net of recovery, $E_k(1-R_k)$, so
\begin{equation}
\label{eq:loss}
L=\sum_{k=1}^{100}E_k(1-R_k)I_k,
\end{equation}
with $\E[L]=\sum_kE_k(1-R_k)\pi_k$ and $\E[L\mid\Theta]=\sum_kE_k(1-R_k)p_k(\Theta)$.
$\E[L]$ is the same under every model of $\Theta$ satisfying $\E[p_k(\Theta)]=\pi_k$; $\E[L\mid\Theta]$ is the loss of an infinitely granular (``large pool'') portfolio with the same systematic risk.

\paragraph{(iii) Standard deviation of $Y_k$.}
By independence of $\epsilon_k$ and $\Theta$, $\Var(Y_k)=\lambda_ka_k'\Sigma a_k+(1-\lambda_k)s_k^2$. Setting $\Var(Y_k)=s_k^2$ gives $\lambda_ks_k^2=\lambda_ka_k'\Sigma a_k$ and, as $\lambda_k>0$,
\begin{equation}
\label{eq:sk}
s_k=\sqrt{a_k'\Sigma a_k}=\sqrt{a_{k1}^2\sigma_1^2+2a_{k1}a_{k2}\rho\sigma_1\sigma_2+a_{k2}^2\sigma_2^2}.
\end{equation}

\paragraph{Two consequences.}
First, $\Corr(Y_j,Y_k)=\sqrt{\lambda_j\lambda_k}\,\Corr(a_j'\Theta,a_k'\Theta)$, so $\lambda_k$ plays the role of the IRB correlation $\rho_k$ (Section~\ref{sec:irb}). Second, with $V_k=a_k'(\Theta-\mu)/s_k$ (mean zero, unit variance),
\begin{equation}
\label{eq:standardised}
Y_k\le d_k\iff\sqrt{\lambda_k}\,V_k+\sqrt{1-\lambda_k}\,\epsilon_k\le d_k^*:=\frac{d_k-\sqrt{\lambda_k}\,a_k'\mu}{s_k}.
\end{equation}
Once $d_k$ is calibrated to $\pi_k$, the loss depends on $\Theta$ only through the \emph{shape} of the law of $(V_1,\dots,V_{100})$, not through the level of means and volatilities. This is why we can impose the one-year $\pi_k$ at the one-week horizon: the horizon enters only through the non-normality of weekly returns, and the comparison with the one-year IRB formula stays like for like.

\subsection*{B. ESTIMATION: QUESTION (iv)}
\paragraph{M1.} Each observed pair $\theta_t$, $t=1,\dots,n$, has probability $1/n$; pairs are kept together, preserving the empirical dependence.

\paragraph{M2.} For an i.i.d.\ Gaussian sample the log-likelihood is
\[
\ell(\mu,\Sigma)=-n\log(2\pi)-\tfrac n2\log|\Sigma|-\tfrac12\textstyle\sum_t(\theta_t-\mu)'\Sigma^{-1}(\theta_t-\mu).
\]
The condition $\partial\ell/\partial\mu=\Sigma^{-1}\sum_t(\theta_t-\mu)=0$ gives $\hat\mu=\bar\theta$. Substituting it, the profile log-likelihood is $-\frac n2\{\log|\Sigma|+\mathrm{tr}(\Sigma^{-1}S)\}$ with $S=n^{-1}\sum_t(\theta_t-\bar\theta)(\theta_t-\bar\theta)'$. Differentiating with respect to $\Sigma^{-1}$, using $\partial\log|\Sigma^{-1}|/\partial\Sigma^{-1}=\Sigma$ and $\partial\,\mathrm{tr}(\Sigma^{-1}S)/\partial\Sigma^{-1}=S$, gives $\hat\Sigma=S$, i.e.
\[
\hat\sigma_i^2=\frac1n\sum_t(\theta_{ti}-\bar\theta_i)^2,\qquad \hat\rho=\frac{\frac1n\sum_t(\theta_{t1}-\bar\theta_1)(\theta_{t2}-\bar\theta_2)}{\hat\sigma_1\hat\sigma_2}.
\]
Standard errors are $\hat\sigma_i/\sqrt n$ for $\hat\mu_i$, $\hat\sigma_i/\sqrt{2n}$ for $\hat\sigma_i$ and $(1-\hat\rho^2)/\sqrt n$ for $\hat\rho$.

\paragraph{M3.} Let $z_i=(\theta_i-\mu_i)/\sigma_i$ and $u_i=\Phi(z_i)$. By Sklar's theorem \citep[Sect.~5.1]{McNeil05} the density of $\Theta$ is the copula density times the marginal densities, $f_\Theta=c_{\nu,\rho}(u_1,u_2)\prod_i\sigma_i^{-1}\varphi(z_i)$. The $t$-copula is $C_{\nu,\rho}(u_1,u_2)=F_{t,2}(t_\nu^{-1}(u_1),t_\nu^{-1}(u_2);\nu,\rho)$; differentiating in $u_1,u_2$ with $\partial t_\nu^{-1}(u)/\partial u=1/f_{t,1}(t_\nu^{-1}(u);\nu)$ gives
\[
c_{\nu,\rho}(u_1,u_2)=\frac{f_{t,2}(x_1,x_2;\nu,\rho)}{f_{t,1}(x_1;\nu)f_{t,1}(x_2;\nu)},\qquad x_i=t_\nu^{-1}(\Phi(z_i)),
\]
with $f_{t,1}$, $f_{t,2}$ the densities of Section~6.2 of the assignment. The log-likelihood of $\psi=(\mu_1,\mu_2,\sigma_1,\sigma_2,\rho,\nu)$ is
\begin{multline*}
\ell(\psi)=\sum_{t=1}^n\Big\{\sum_{i=1}^2\big[-\log\sigma_i-\tfrac12\log(2\pi)-\tfrac12z_{ti}^2\big]\\
+\log f_{t,2}(x_{t1},x_{t2};\nu,\rho)-\sum_{i=1}^2\log f_{t,1}(x_{ti};\nu)\Big\}.
\end{multline*}
The margins enter the copula term through $x_{ti}$, so there is no closed form. We first estimate the margins by their Gaussian likelihood and $(\rho,\nu)$ from the copula term (inference functions for margins, \citealp[Sect.~5.5]{McNeil05}), then maximise $\ell$ jointly in all six parameters (L-BFGS-B, $|\rho|<1$, $2<\nu\le200$, optimising $\log\nu$). Standard errors come from the inverse observed information. M2 is the limit $\nu\to\infty$, on the boundary of the parameter space, so the likelihood-ratio statistic follows a $\frac12\chi^2_0+\frac12\chi^2_1$ mixture \citep{SelfLiang87}. The coefficient of tail dependence is $\lambda_L=2\,t_{\nu+1}\big(-\sqrt{(\nu+1)(1-\rho)/(1+\rho)}\big)$ \citep[Example~5.33]{McNeil05}, against $0$ for the Gaussian copula; it equals \val{q4.M3.tail_dep|f3} at the estimates. The Pearson correlation of $\Theta$ under M3, needed in (\ref{eq:sk}), is computed by numerical integration.

\subsection*{C. SIMULATION AND THRESHOLDS: QUESTIONS (v)--(vii)}
\paragraph{Factor draws.} M1: draw $t$ uniformly on $\{1,\dots,n\}$ and set $\Theta=\theta_t$. M2: $\Theta=\hat\mu+A'Z$, $Z\sim N(0,I_2)$, with $A$ the upper-triangular Cholesky factor, $A'A=\hat\Sigma$. M3: $Z\sim N(0,P)$ with correlation $\hat\rho$ and $W\sim\chi^2_{\hat\nu}$ independent; $X=\sqrt{\hat\nu/W}\,Z$ is bivariate $t$, $U_i=t_{\hat\nu}(X_i)$ has the $t$-copula, and $\Theta_i=\hat\mu_i+\hat\sigma_i\Phi^{-1}(U_i)$ has the required Gaussian margin by the quantile transformation. Each scenario adds $100$ independent $\epsilon_k$ and applies (\ref{eq:Yk}) and (\ref{eq:loss}).

\paragraph{Common random numbers.} Every scenario uses the same primitive draws for all models: the Gaussian pair $Z$ enters M2 and M3, the mixing variable of M3 is $W=F^{-1}_{\chi^2_{\hat\nu}}(U_W)$ for a uniform $U_W$, M1 uses the week $\lceil nU_B\rceil$ for a uniform $U_B$, and the shocks $\epsilon_k$ are shared. All generators are seeded.

\paragraph{Empirical-quantile threshold.} With $N$ draws of $Y_k$, the threshold is the generalised inverse of the empirical distribution at $\pi_k$:
\[
\hat d_k=\inf\Big\{y:\tfrac1N\textstyle\sum_j\ind_{\{Y_k^{(j)}\le y\}}\ge\pi_k\Big\}=Y_{k,(\lceil N\pi_k\rceil)} .
\]
$N\pi_k$ is an integer for all ratings, so each firm defaults in exactly a share $\pi_k$ of the scenarios. For \val{q5.n_names_k_le_7|int} names this rests on seven or fewer draws.

\paragraph{Reference threshold.} We solve $\E[p_k(\Theta;d_k)]=\pi_k$ on the model, integrating $\epsilon_k$ analytically through (\ref{eq:condpd}). M2: $d_k=\sqrt{\lambda_k}a_k'\mu+s_k\Phi^{-1}(\pi_k)$. M1: $d_k$ is the root of $n^{-1}\sum_t p_k(\theta_t;d_k)=\pi_k$. M3: the expectation is a two-dimensional integral against the density of $(G_1,G_2)$, $c_{\nu,\rho}(\Phi(g_1),\Phi(g_2))\varphi(g_1)\varphi(g_2)$, computed with the trapezoidal rule (step $0.1$ on $[-8.5,8.5]^2$; a five times finer grid changes the implied probabilities by less than $10^{-5}$). Roots are found by safeguarded Newton iteration.

\paragraph{Reference run and rolling windows.} The reference run has $2{,}000{,}000$ scenarios per model in $200$ blocks of $10{,}000$ with the reference thresholds; the blocks are independent replications of the baseline. In the rolling analysis each window uses overlapping five-day returns (a pseudo-likelihood), quadrature step $0.2$ and $50{,}000$ scenarios with the same random numbers at every date. The placebo band comes from M3 with the parameters of M2 and $\nu=200$, run as one window for $200$ seeds.

\subsection*{D. VALUE-AT-RISK AND EXPECTED SHORTFALL}
With $X=-L$ \citep[Ch.~1]{DeGiorgi26},
\begin{align*}
\VaR_\alpha(X)&=-\inf\{x:\P[X\le x]\ge1-\alpha\},\\
\ES_\alpha(X)&=\tfrac{1}{1-\alpha}\Big\{\E\big[-X\ind_{\{X\le-\VaR_\alpha\}}\big]-\VaR_\alpha\big(\P[X\le-\VaR_\alpha]-(1-\alpha)\big)\Big\}.
\end{align*}
The second ES term matters because $L$ is discrete. For ordered simulated losses $\ell_{(1)}\ge\dots\ge\ell_{(N)}$, $\widehat\VaR_\alpha=\ell_{(k^*)}$ with $k^*=\lceil N(1-\alpha)\rceil$ (the tenth-largest loss for $N=10{,}000$), and
\[
\widehat\ES_\alpha=\frac{1}{N(1-\alpha)}\Big\{\sum_{j:\ell_j\ge\widehat\VaR_\alpha}\ell_j-\widehat\VaR_\alpha\big(\#\{j:\ell_j\ge\widehat\VaR_\alpha\}-N(1-\alpha)\big)\Big\},
\]
the average of the ten largest losses, corrected for ties.

\subsection*{E. CLIMATE OVERLAY: QUESTION (xi)}
\paragraph{Loss with storms.} With $m_k$ the worst damage ratio of the year at $k$ and $\phi_k=\phi$,
\[
L^{\mathrm{clim}}=\sum_kE_k\big(1-R_k(1-\phi m_k)\big)I_k=L+\phi A,
\]
with $A=\sum_kE_kR_km_kI_k\ge0$, so EL, VaR and ES can only increase.

\paragraph{Exact law of the worst event.} Order the events damaging $k$ by decreasing ratio, $m_{k,(1)}\ge m_{k,(2)}\ge\dots$, with frequencies $f_{(j)}$. The worst event of the year is the $j$-th if none of the worse ones occurs and the $j$-th does:
\[
\P[m_k=m_{k,(j)}]=e^{-(f_{(1)}+\dots+f_{(j-1)})}\big(1-e^{-f_{(j)}}\big),\quad \P[m_k=0]=e^{-\sum_jf_{(j)}} .
\]
By independence, $\E[L^{\mathrm{clim}}]-\E[L]=\phi\sum_kE_kR_k\pi_k\E[m_k]$ exactly.

\paragraph{First-order tail effect.} For a small perturbation \citep{Gourieroux00}, $\Delta\VaR_\alpha\approx\phi\,\E[A\mid L=\VaR_\alpha]$ and $\Delta\ES_\alpha\approx\phi\,\E[A\mid L\ge\VaR_\alpha]$, with $\E[A\mid L]=\sum_kE_kR_k\E[m_k]\P[I_k=1\mid L]$ under independence. Only the expected damage at the locations of the defaulting names matters to first order, which is why concentration is irrelevant under independence. For the reference run we use the equivalent sampler $M\sim\mathrm{Poisson}(\sum_ef_e)$ events with identities drawn with probabilities $f_e/\sum f$ (Poisson superposition).

\paragraph{Link to \citet{Pozdyshev25}.} In their model an event scales assets by $e^{-\alpha}$, giving $\mathrm{LGD}_1=\mathrm{LGD}_0+(1-e^{-\alpha})(1-\mathrm{LGD}_0)$; with $1-R_k(e)=\mathrm{LGD}_{0,k}+\phi\MDR_k(e)(1-\mathrm{LGD}_{0,k})$ our overlay is this LGD channel with $1-e^{-\alpha}=\phi\MDR$, and their add-on reproduces our expected-loss effect to machine precision. For the PD channel the worst damage shifts the default threshold by $\alpha/\sigma_A$ standard deviations, $I_k=\ind\{Y_k\le d_k+s_k\alpha/\sigma_A\}$ with $\alpha=-\log(1-\phi m_k)$ and $\sigma_A=25\%$; the default probability becomes the $q$-deformed mixture $\sum_jp_j\Phi(\Phi^{-1}(\pi_k)+\alpha_j/\sigma_A)$. The defaults added in the tail are replaced by their expectation given $\Theta$ to obtain a first-order $\Delta\ES$. The gust stress multiplies every gust speed by $1.05$, $1.10$ or $1.20$ and re-evaluates the impact function.

\subsection*{F. CODE AND CORRESPONDENCE WITH THE ASSIGNMENT}
\sloppy The Python code is attached. The package \texttt{src/qrm} has one module per component (\texttt{data}, \texttt{factor\_models}, \texttt{credit\_model}, \texttt{risk\_measures}, \texttt{rolling}, \texttt{basel\_irb}, \texttt{climada\_io}, \texttt{climate}, \texttt{climate\_channels}); the numbered scripts reproduce every table, figure and number of this report, and \texttt{run\_all.py} runs them in sequence. Only script \texttt{06} needs CLIMADA; it stores the storm catalogue, so the rest runs without it. All random draws are seeded, and a fresh run reproduces the reported numbers exactly. Table~\ref{tab:map} maps each question to this report.

\begin{table}[!htb]
\caption{Where each question of the assignment is answered.}
\label{tab:map}
\centering
{\small\setlength{\tabcolsep}{4pt}
\begin{tabular}{@{}lp{0.5\textwidth}l@{}}
\toprule
Item & Content & Where \\
\midrule
(i) & Conditional default probability; law of $I_k$ given $\Theta$ & App.~A, eq.~(\ref{eq:condpd}) \\
(ii) & Portfolio loss & App.~A, eq.~(\ref{eq:loss}) \\
(iii) & Standard deviation of $Y_k$ & App.~A, eq.~(\ref{eq:sk}) \\
(iv) & ML estimation of M2 and M3 & Sect.~\ref{sec:setting}, Table~\ref{tab:mle}; App.~B \\
(v)--(vi) & Simulation of $Y_k$; thresholds $d_k$ & Sect.~\ref{sec:method}; App.~C \\
(vii) & Portfolio loss distribution & Sect.~\ref{sec:method}, Fig.~\ref{fig:hist} \\
(viii) & $99.9\%$ VaR and ES; comparison & Sect.~\ref{sec:results}, Tables~\ref{tab:var}--\ref{tab:varref}; App.~D \\
(ix) & VaR and ES on rolling 500-day windows & Sect.~\ref{sec:rolling}, Fig.~\ref{fig:rolling} \\
(x) & $\rho_k$, $K_k$, meaning of $\mathrm{MA}(\pi_k,1)$, $K$, RWA & Sect.~\ref{sec:irb}, Table~\ref{tab:irb} \\
(x) & IRB capital vs.\ UL; regulator's view & Sect.~\ref{sec:irb}, Table~\ref{tab:gap} \\
(xi) & Interpretation of $\phi_k=0$ and $\phi_k=1$ & Sect.~\ref{sec:climate} \\
(xi) & Storm overlay: VaR and ES vs.\ (viii) & Sect.~\ref{sec:climate}, Table~\ref{tab:clim}; App.~E \\
(xi) & Geographically diversified portfolio; canton level & Sect.~\ref{sec:climate}, Table~\ref{tab:climvar} \\
\bottomrule
\end{tabular}}
\end{table}


% =====================================================================
% bibliography.tex  --  natbib author-year, format of roffin.cls.
% Course materials first; external sources only where a specific result
% is used (the place is indicated in the text).
% =====================================================================

\begin{thebibliography}{}

\bibitem[{Aznar-Siguan and Bresch(2019)}]{AznarSiguan19}
Aznar-Siguan, G. and Bresch, D.~N. (2019)
CLIMADA v1: a global weather and climate risk assessment platform,
\textit{Geoscientific Model Development} \textbf{12}, 3085--3097.

\bibitem[{Basel Committee on Banking Supervision(2005)}]{BCBS05}
Basel Committee on Banking Supervision (2005)
\textit{An Explanatory Note on the Basel II IRB Risk Weight Functions}.
Basel: Bank for International Settlements.

\bibitem[{De~Giorgi(2026)}]{DeGiorgi26}
De~Giorgi, E. (2026)
\textit{Quantitative Risk Management}, Lecture Notes,
University of St.~Gallen.

\bibitem[{De~Giorgi and Filkov\'a(2026)}]{DeGiorgiFilkova26}
De~Giorgi, E. and Filkov\'a, Z. (2026)
\textit{Quantitative Risk Management: Assignment},
University of St.~Gallen, 12 October 2026.

\bibitem[{Gordy(2003)}]{Gordy03}
Gordy, M.~B. (2003)
A risk-factor model foundation for ratings-based bank capital rules,
\textit{Journal of Financial Intermediation} \textbf{12}, 199--232.

\bibitem[{Gouri\'eroux et~al.(2000)}]{Gourieroux00}
Gouri\'eroux, C., Laurent, J.-P. and Scaillet, O. (2000)
Sensitivity analysis of Values at Risk,
\textit{Journal of Empirical Finance} \textbf{7}, 225--245.

\bibitem[{Hauser et~al.(2026)}]{HauserEtAl26}
Hauser, M., Noyelle, R. and Juhel, S. (2026)
\textit{Climate risk modelling with CLIMADA: introduction and mandatory exercises},
course material, ETH Zurich, Spring Term 2026.

\bibitem[{McNeil et~al.(2005)}]{McNeil05}
McNeil, A.~J., Frey, R. and Embrechts, P. (2005)
\textit{Quantitative Risk Management: Concepts, Techniques and Tools}.
Princeton: Princeton University Press.

\bibitem[{Self and Liang(1987)}]{SelfLiang87}
Self, S.~G. and Liang, K.-Y. (1987)
Asymptotic properties of maximum likelihood estimators and likelihood ratio tests under nonstandard conditions,
\textit{Journal of the American Statistical Association} \textbf{82}, 605--610.

\bibitem[{Vasicek(2002)}]{Vasicek02}
Vasicek, O. (2002)
The distribution of loan portfolio value,
\textit{Risk} \textbf{15}(12), 160--162.

\bibitem[{Welker et~al.(2021)}]{Welker21}
Welker, C., R\"o\"osli, T. and Bresch, D.~N. (2021)
Comparing an insurer's perspective on building damages with modelled damages from pan-European winter windstorm event sets: a case study from Zurich, Switzerland,
\textit{Natural Hazards and Earth System Sciences} \textbf{21}, 279--299.

\end{thebibliography}


\end{document}
